\chapter{梁惠王章句上}
\kaishi*

%梁惠王章句上
\jiazhu[1]{梁惠王者，魏惠王也。魏，國名；惠，謚也；王，號也。時天下有七王，皆僭號者也。猶春秋之時，吳楚之君稱王也。魏惠王居於大梁，故號曰梁王。聖人及大賢有道德者，王公侯伯及卿大夫咸願以爲師。孔子時，諸侯問疑質禮，若弟子之問師也。魯衞之君皆尊事焉，故論語或以弟子名篇，而有衞靈公、季氏之篇。孟子亦以大儒爲諸侯所師，是以梁惠王、滕文公題篇與公孫丑等爲一例也。}

\section{孟子見梁惠王}

孟子見梁惠王。\jiazhu[1]{孟子適梁，魏惠王禮請孟子見之。}

王曰：「叟！不遠千里而來，亦將有以利吾國乎？」\jiazhu[1]{曰，辭也。叟，長老之稱也，猶父也。孟子去齊，老而之魏，故王尊禮之曰父。不遠千里之路而來至此，亦將有可以爲寡人興利除害也。}

孟子對曰：「王何必曰利？亦有仁義而已矣。」\jiazhu[1]{孟子知王欲以富國彊兵爲利，故曰王何必以利爲名乎？亦惟有仁義之道者可以爲名。以利爲名，則有不利之患矣。因爲王陳之。}

「王曰：『何以利吾國？』大夫曰：『何以利吾家？』士庶人曰：『何以利吾身？』上下交征利，而國危矣。」\jiazhu[1]{征，取也。從王至庶人，故言上下交爭各欲利其身，必至於篡弑，則國危亡矣。論語曰：「放於利而行，多怨。」故不欲使王以利爲名也。又言交爲俱也。}

「萬乘之國弑其君者，必千乘之家；千乘之國殺其君者，必百乘之家。」\jiazhu[1]{萬乘，兵車萬乘，謂天子也。千乘，兵車千乘，謂諸侯也。夷羿之弑夏后，是以千乘取萬乘也。天子建國，諸侯立家。百乘之家，謂大國之卿，食菜邑有兵車百乘之賦者也。若齊崔、衞寗、晉六卿等，是以其終亦皆弑其君。此以百乘取千乘也。上千乘當言國而言家者，諸侯以國爲家，亦以避萬乘稱國，故稱家，君臣上下之辭。}

「萬取千焉，千取百焉，不爲不多矣。」\jiazhu[1]{周制，君十卿禄。君食萬鍾，臣食千鍾，亦多矣，不爲不多矣。}

「苟爲後義而先利，不奪不饜。」\jiazhu[1]{苟，誠也。誠令大臣皆後仁義而先自利，則不篡奪君位，不足自饜飽其欲。}

「未有仁而遺其親者也，未有義而後其君者也。」\jiazhu[1]{仁者親親，義者尊尊。人無行仁而遺棄其親，行義而忽後其君者。}

「王亦曰仁義而已矣，何必曰利？」\jiazhu[1]{孟子復申此者，重嗟歎其禍。} \jiazhu[1]{}\par \jiazhu[1]{章指言：治國之道，明當以仁義爲名，然後上下和親，君臣集穆，天經地義，不易之道，故以建篇立始也。}

\section{孟子見梁惠王於沼上}

孟子見梁惠王。王立於沼上，顧鴻鴈麋鹿，曰：「賢者亦樂此乎？」\jiazhu[1]{沼，池也。王好廣苑囿，大池沼，與孟子遊觀，顧視禽獸之衆多，心以爲娛樂，夸咤孟子曰：「賢者亦樂此乎？」}

孟子對曰：「賢者而後樂此，不賢者雖有此，不樂也。」\jiazhu[1]{惟有賢者然後乃得樂此耳。謂脩堯舜之道，國家安寧，故得有此以爲樂也。不賢之人，亡國破家，雖有此，當爲人所奪，故不得以爲樂也。}

「《詩》云：『經始靈臺，經之營之。庶民攻之，不日成之。經始勿亟，庶民子來。王在靈囿，麀鹿攸伏。麀鹿濯濯，白鳥鶴鶴。王在靈沼，於牣魚躍。』」\jiazhu[1]{詩大雅靈臺之篇也。言文王始經營規度此臺，衆民竝來治作之，不與期日，自來成之也。言文王戒固勿亟疾，而衆民自來，如子趣父事，自來成之也。麀鹿，牝鹿也。言文王在囿中，麀鹿懷任，安其所而伏，不驚動也。獸肥飽則濯濯，鳥肥飽則鶴鶴而澤好。文王在池沼，魚乃跳躍喜樂，言其德及鳥獸魚鼈也。}

「文王以民力爲臺爲沼，而民歡樂之，謂其臺曰靈臺，謂其沼曰靈沼，樂其有麋鹿魚鼈。」\jiazhu[1]{文王雖以民力築臺鑿池，民由歡樂之，謂其臺沼若神靈之所爲，欲使其多禽獸以養文王者也。}

「古之人與民偕樂，故能樂也。」\jiazhu[1]{言古之賢君與民共同其所樂，故能樂之。}

「《湯誓》曰：『時日害喪，予及女皆亡。』」\jiazhu[1]{湯誓，尚書篇名也。時，是也。時乙卯日也。害，大也。言桀爲無道，百姓皆欲與湯共伐之。湯臨士衆而誓之言：「是日桀當大喪亡，我與汝俱往亡之。」}

「民欲與之皆亡，雖有臺池鳥獸，豈能獨樂哉？」\jiazhu[1]{孟子說詩書之義以感喻王。言民皆欲與湯共亡桀，雖有臺池禽獸，何能復獨樂之哉？復申明上言「不賢者雖有此不樂」也。}\par \jiazhu[1]{章指言：聖王之德，與民共樂，恩及鳥獸，則忻戴其上，太平化興；無道之君，衆怨神怒，則國滅祀絶，不得保守其所樂也。}

\section{梁惠王問民不加多}

梁惠王曰：「寡人之於國也，盡心焉耳矣。」\jiazhu[1]{王侯自稱孤寡。言寡人於治國之政，盡心欲利百姓。焉耳者，懇至之辭。}

「河內凶，則移其民於河東，移其粟於河內。河東凶亦然。」\jiazhu[1]{言凶年以此救民也。魏舊在河東，後爲強國，兼得河內也。}

「察鄰國之政，無如寡人之用心者。」\jiazhu[1]{言鄰國之君用心憂民無如已也。}

「鄰國之民不加少，寡人之民不加多，何也？」\jiazhu[1]{王自怪爲政有此惠而民人不增多於鄰國者，何也？}

孟子對曰：「王好戰，請以戰喻。」\jiazhu[1]{因王好戰，故以戰事喻解王意。}

「塡然鼓之，兵刃既接，棄甲曳兵而走。或百步而後止，或五十步而後止。以五十步笑百步，則何如？」\jiazhu[1]{塡，鼓音也。兵以鼓進，以金退。孟子問王曰：「今有戰者，兵刃巳交，其負者棄甲曳兵而走，五十步而止，足以笑百步止者不？」}

曰：「不可，直不百步耳，是亦走也。」\jiazhu[1]{王曰：「不足以相笑也。是人俱走，直事不百步耳。」}

曰：「王如知此，則無望民之多於鄰國也。」\jiazhu[1]{孟子曰：「王如此不足以相笑。王之政猶此也。王雖有移民轉穀之善政，其好戰殘民與鄰國同，而獨望民之多，何異於以五十步笑百步者乎？」}

「不違農時，穀不可勝食也。」\jiazhu[1]{從此巳下，爲王陳王道也。使民得三時務農，不違奪其要時，則五穀饒穰，不可勝食。}

「數罟不入洿池，魚鼈不可勝食也。」\jiazhu[1]{數罟，密網也。密細之網所以捕小魚鼈者也，故禁之不得用。魚不滿尺，不得食。}

「斧斤以時入山林，材木不可勝用也。」\jiazhu[1]{時謂草木零落之時，使林木茂暢，故有餘。}

「穀與魚鼈不可勝食，材木不可勝用，是使民養生喪死無憾也。」\jiazhu[1]{憾，恨也。民所用者足，故無恨。}

「養生喪死無憾，王道之始也。」\jiazhu[1]{王道先得民心，民心無恨，故言王道之始。}

「五畒之宅，樹之以桑，五十者可以衣帛矣。」\jiazhu[1]{廬井邑居各二畒半以爲宅，冬入保城二畒半，故爲五畒也。樹桑牆下，古者年五十乃衣帛矣。}

「雞豚狗彘之畜，無失其時，七十者可以食肉矣。」\jiazhu[1]{言孕字不失時也。七十不食肉不飽。}

「百畒之田，勿奪其時，數口之家可以無飢矣。」\jiazhu[1]{一夫一婦耕耨百畒。百畒之田不可以傜役奪其時功，則家給人足。農夫上中下所食多少各有差，故揔言數口之家也。}

「謹庠序之敎，申之以孝悌之義，頒白者不負戴於道路矣。」\jiazhu[1]{庠序者，敎化之宫也。殷曰序，周曰庠。謹脩敎化，申重孝悌之義。頒者，斑也，頭半白斑斑者也。壯者代老，心各安之，故斑白者不負戴也。}

「七十者衣帛食肉，黎民不飢不寒，然而不王者，未之有也。」\jiazhu[1]{言百姓老稚温飽，禮義脩行，積之可以致王也。孟子欲以風王，何不行此，可以王天下，有率土之民，何但望民多於鄰國。}

「狗彘食人食而不知檢，塗有餓莩而不知發。」\jiazhu[1]{言人君但養犬彘使食人食，不知以法度檢斂也。塗，道也。餓死者曰莩。詩曰：「莩有梅。」莩，零落也。道路之旁有餓死者，不知發倉廩以用振救之也。}

「人死，則曰：『非我也，歲也。』是何異於刺人而殺之，曰：『非我也，兵也。』」\jiazhu[1]{人死謂餓疫死者也。王政使然，而曰非我殺之，歲殺之也。此何以異於用兵殺人，而曰非我也，兵自殺之也。}

「王無罪歲，斯天下之民至焉。」\jiazhu[1]{戒王無歸罪於歲，責已而改行，則天下之民皆可致也。}\par \jiazhu[1]{章指言：王化之本，在於使民養生喪死之用僃足，然後導之以禮義，責已矜窮，則斯民集矣。}

\section{梁惠王問承教}

梁惠王曰：「寡人願安承敎。」\jiazhu[1]{願安意承受孟子之敎令。}

孟子對曰：「殺人以梃與刃，有以異乎？」\jiazhu[1]{梃，杖也。}

曰：「無以異也。」\jiazhu[1]{王曰：「杖刃殺人，無以異也。」}

「以刃與政，有以異乎？」\jiazhu[1]{孟子欲以次喻王。}

曰：「無以異也。」\jiazhu[1]{王復曰：「政殺人，無以異也。」}

曰：「庖有肥肉，廏有肥馬，民有飢色，野有餓莩，此率獸而食人也。」\jiazhu[1]{孟子言人君如此爲率禽獸以食人也。}

「獸相食，且人惡之；爲民父母，行政不免於率獸而食人，惡在其爲民父母也？」\jiazhu[1]{虎狼食禽獸，人猶尚惡視之。牧民爲政乃率禽獸食人，安在其爲民父母之道也。}

「仲尼曰：『始作俑者，其無後乎！』爲其象人而用之也。如之何其使斯民飢而死也？」\jiazhu[1]{俑，偶人也，用之送死。仲尼重人類，謂秦穆公時以三良殉葬，本由有作俑者也。夫惡其始造，故曰此人其無後嗣乎！如之何其使此民飢而死邪？孟子陳此以敎王愛民。}\par \jiazhu[1]{章指言：王者爲政之道，生民爲首，以政殺人，人君之咎，猶以白刃疾之甚也。}

\section{梁惠王問國恥}

梁惠王曰：「晉國，天下莫強焉，叟之所知也。」\jiazhu[1]{韓、魏、趙本晉六卿，當此時號三晉，故惠王言晉國天下強也。}

「及寡人之身，東敗於齊，長子死焉；西喪地於秦七百里；南辱於楚。寡人恥之，願比死者壹洒之，如之何則可？」\jiazhu[1]{王念有此三恥，求䇿謀於孟子。}

孟子對曰：「地方百里而可以王。」\jiazhu[1]{言古聖人以百里之地以致王天下，謂文王也。}

「王如施仁政於民，省刑罰，薄稅斂，深耕易耨，壯者以暇日脩其孝悌忠信，入以事其父兄，出以事其長上，可使制梃以撻秦楚之堅甲利兵矣。」\jiazhu[1]{易耨，芸苗令簡易也。制，作也。王如行此政，可使國人作杖以捶敵國堅甲利兵，何患恥之不雪也。}

「彼奪其民時，使不得耕耨以養其父母。父母凍餓，兄弟妻子離散。彼陷溺其民，王往而征之，夫誰與王敵？」\jiazhu[1]{彼謂齊秦楚也。彼困其民，願王往征之也。彼失民心，民不爲用，夫誰與共禦王之師，爲王敵乎？}

「故曰：『仁者無敵。』王請勿疑。」\jiazhu[1]{鄰國暴虐，已脩仁政，則無敵矣。王請行之，勿有疑也。}\par \jiazhu[1]{章指言：以百里行仁，天下歸之；以政傷民，民樂其亡；以梃服強，仁與不仁也。}

\section{孟子見梁襄王}

孟子見梁襄王。出，語人曰：「望之不似人君。」\jiazhu[1]{襄，謚也，魏之嗣王也。望之無儼然之威儀也。}

「就之而不見所畏焉。」\jiazhu[1]{就與之言，無人君操秉之威，知其不足畏。}

「卒然問曰：『天下惡乎定？』」\jiazhu[1]{卒暴問事，不由其次也。問天下安所定，言誰能定之。}

「吾對曰：『定于一。』」\jiazhu[1]{孟子謂仁政爲一也。}

「『孰能一之？』」\jiazhu[1]{言孰能一之者。}

「對曰：『不嗜殺人者能一之。』」\jiazhu[1]{嗜，猶甘也。言令諸侯有不甘樂殺人者，則能一之。}

「『孰能與之？』」\jiazhu[1]{王言誰能與不嗜殺人者乎。}

「對曰：『天下莫不與也。』」\jiazhu[1]{孟子曰：時人皆苦虐政，如有行仁，天下莫不與之。}

「王知夫苗乎？七八月之閒旱，則苗槁矣。天油然作雲，沛然下雨，則苗浡然興之矣。其如是，孰能禦之？」\jiazhu[1]{以苗生喻人象也。周七八月，夏之五六月。油然，興雲之貌。沛然，下雨以潤槁苗，則浡然巳盛，孰能止之。}

「今夫天下之人牧，未有不嗜殺人者也。如有不嗜殺人者，則天下之民皆引領而望之矣。誠如是也，民歸之，由水之就下，沛然誰能禦之？」\jiazhu[1]{今天下牧民之君，誠能行此仁政，民皆延頸望欲歸之，如水就下，沛然而來，誰能止之。}\par \jiazhu[1]{章指言：定天下者一道而已，不貪殺人，人則歸之。是故文王視民如傷，此之謂也。}

\section{齊宣王問桓文之事}

齊宣王問曰：「齊桓、晉文之事，可得聞乎？」\jiazhu[1]{宣，謚也。宣王問孟子欲庶幾齊桓公小白、晉文公重耳。孟子兾得行道，故仕於齊，不用而去，乃適於梁。建篇先梁者，欲以仁義首篇，因言魏事，章次相從，然後道齊也。}

孟子對曰：「仲尼之徒無道桓、文之事者，是以後世無傳焉，臣未之聞也。」\jiazhu[1]{孔子之門徒，頌述宓戲以來至文武周公之法制耳，雖及五霸，心賤薄之。是以儒家後世無欲傳道之者，故曰臣未之聞也。}

「無以，則王乎？」\jiazhu[1]{旣不論三皇五帝，殊無所問，則尚當問王道耳。不欲使王問霸事也。}

曰：「德何如，則可以王矣？」\jiazhu[1]{王曰：德行當何如而可以王乎。}

曰：「保民而王，莫之能禦也。」\jiazhu[1]{保，安也。禦，止也。言安民則惠，黎民懷之，若此以王，無能止也。}

曰：「若寡人者，可以保民乎哉？」\jiazhu[1]{王自恐德不足以安民，故問之。}

曰：「可。」\jiazhu[1]{孟子以爲如王之性，可以安民也。}

曰：「何由知吾可也？」\jiazhu[1]{王問孟子何以知吾可以安民。}

曰：「臣聞之胡齕曰，王坐於堂上，有牽牛而過堂下者，王見之，曰：『牛何之？』對曰：『將以釁鍾。』王曰：『舍之！吾不忍其觳觫，若無罪而就死地。』對曰：『然則廢釁鍾與？』曰：『何可廢也？以羊易之！』不識有諸？」\jiazhu[1]{胡齕，王左右近臣也。觳觫，牛當到死地處恐貌。新鑄鍾，殺牲以血塗其釁郄，因以祭之曰釁。周禮大祝曰：「墮釁逆牲逆尸，令鍾鼓。」天府上春釁寶鍾及寶器。孟子曰：臣受胡齕言王甞有此仁，不知誠有之否。}

曰：「有之。」\jiazhu[1]{王曰有之。}

曰：「是心足以王矣。百姓皆以王爲愛也，臣固知王之不忍也。」\jiazhu[1]{愛，嗇也。孟子曰：王推是仁心，足以至於王道。然百姓皆謂王嗇愛其財，臣知王見牛恐懼，不欲趨死，不忍故易之也。}

王曰：「然，誠有百姓者。齊國雖褊小，吾何愛一牛？即不忍其觳觫，若無罪而就死地，故以羊易之也。」\jiazhu[1]{王曰：亦誠有百姓所言者矣。吾國雖小，豈愛惜一牛之財費哉？即見其牛哀之，釁鍾又不可廢，故易之以羊耳。}

曰：「王無異於百姓之以王爲愛也。以小易大，彼惡知之？王若隱其無罪而就死地，則牛羊何擇焉？」\jiazhu[1]{異，怪也。隱，痛也。孟子言無怪百姓之謂王愛財也，見王以小易大故也。王如痛其無罪，羊亦無罪，何爲獨釋牛而取羊。}

王笑曰：「是誠何心哉？我非愛其財而易之以羊也。宜乎百姓之謂我愛也。」\jiazhu[1]{王自笑心不然，而不能自免爲百姓所非，乃責已之以小易大，故曰宜乎其罪我也。}

曰：「無傷也，是乃仁術也，見牛未見羊也。君子之於禽獸也，見其生，不忍見其死；聞其聲，不忍食其肉。是以君子遠庖廚也。」\jiazhu[1]{孟子解王自責之心曰：無傷於仁，是乃王爲仁之道也。時未見羊，羊之爲牲次於牛，故用之耳。是以君子遠庖廚，不欲見其生食其肉也。}

王說曰：「《詩》云：『他人有心，予忖度之。』夫子之謂也。夫我乃行之，反而求之，不得吾心。夫子言之，於我心有戚戚焉。此心之所以合於王者，何也？」\jiazhu[1]{詩小雅巧言之篇也。王喜恱，因稱是詩以嗟歎孟子忖度知已心。戚戚然，心有動也。寡人雖有是心，何能足以王也。}

曰：「有復於王者曰：『吾力足以舉百鈞，而不足以舉一羽；明足以察秋豪之末，而不見輿薪。』則王許之乎？」\jiazhu[1]{復，白也。許，信也。人有白王如此，王信之乎？百鈞，三千斤也。}

曰：「否。」\jiazhu[1]{王曰我不信也。}

「今恩足以及禽獸，而功不至於百姓者，獨何與？然則一羽之不舉，爲不用力焉；輿薪之不見，爲不用明焉；百姓之不見保，爲不用恩焉。故王之不王，不爲也，非不能也。」\jiazhu[1]{孟子言王恩及禽獸而不安百姓，若不用力、不用明者也。不爲耳，非不能也。}

曰：「不爲者與不能者之形何以異？」\jiazhu[1]{王問其狀何以異也。}

曰：「挾大山以超北海，語人曰『我不能』，是誠不能也。爲長者折枝，語人曰『我不能』，是不爲也，非不能也。故王之不王，非挾大山以超北海之類也；王之不王，是折枝之類也。」\jiazhu[1]{孟子爲王陳爲與不爲之形若是。王則不折枝之類也。折枝，按摩折手節解罷枝也。少者恥見役，故不爲耳，非不能也。大山、北海皆近齊，故以爲喻也。}

「老吾老，以及人之老；幼吾幼，以及人之幼。天下可運於掌。」\jiazhu[1]{老猶敬也，幼猶愛也。敬我之老，亦敬人之老；愛我之幼，亦愛人之幼。推此心以惠民，天下可轉之掌上，言易也。}

「《詩》云：『刑于寡妻，至于兄弟，以御于家邦。』言舉斯心加諸彼而已。」\jiazhu[1]{詩大雅思齊之篇也。刑，正也。寡，少也。言文王正已適妻，則八妾從，以及兄弟。御，享也，享天下國家之福。但舉已心加於人耳。}

「故推恩足以保四海，不推恩無以保妻子。古之人所以大過人者，無他焉，善推其所爲而已矣。」\jiazhu[1]{大過人者，大有爲之君也。善推其心所好惡，以安四海也。}

「今恩足以及禽獸，而功不至於百姓者，獨何與？」\jiazhu[1]{復申此言，非王不能，不爲之耳。}

「權，然後知輕重；度，然後知長短。物皆然，心爲甚。王請度之！」\jiazhu[1]{權，銓衡也，可以稱輕重也。度，丈尺也，可以量長短也。凡物皆當稱度乃可知。心當行之乃爲仁，心比於物尤當爲之甚者也。欲使王度心如度物也。}

「抑王興甲兵，危士臣，搆怨於諸侯，然後快於心與？」\jiazhu[1]{抑，辭也。孟子問王抑亦如是乃快邪。}

王曰：「否，吾何快於是？將以求吾所大欲也。」\jiazhu[1]{王言不然，我不快是也。將欲以求我心所大欲者耳。}

曰：「王之所大欲可得聞與？」\jiazhu[1]{孟子雖心知王意，而故問者，欲令王自道，縁以陳之。}

王笑而不言。\jiazhu[1]{王意大而不敢正言。}

曰：「爲肥甘不足於口與？輕煖不足於體與？抑爲采色不足視於目與？聲音不足聽於耳與？便嬖不足使令於前與？王之諸臣皆足以供之，而王豈爲是哉？」\jiazhu[1]{孟子復問此五者，欲以致王所欲也。故發異端以問也。}

曰：「否，吾不爲是也。」\jiazhu[1]{王言我不爲是也。}

曰：「然則王之所大欲可知已。欲辟土地，朝秦楚，莅中國而撫四夷也。」\jiazhu[1]{莅，臨也。言王意欲庶幾王者莅臨中國而安四夷者也。}

「以若所爲，求若所欲，猶緣木而求魚也。」\jiazhu[1]{若，順也。順嚮者所爲謂搆兵諸侯之事，求順今之所欲莅中國之願。其不可得如緣喬木而求生魚也。}

王曰：「若是其甚與？」\jiazhu[1]{王謂比之緣木求魚爲大甚。}

曰：「殆有甚焉。緣木求魚，雖不得魚，無後災；以若所爲，求若所欲，盡心力而爲之，後必有災。」\jiazhu[1]{孟子言盡心戰鬪，必有殘民破國之災，故曰殆有甚於緣木求魚者也。}

曰：「可得聞與？」\jiazhu[1]{王欲知其害也。}

曰：「鄒人與楚人戰，則王以爲孰勝？」\jiazhu[1]{言鄒小楚大也。}

曰：「楚人勝。」\jiazhu[1]{王曰楚人勝也。}

曰：「然則小固不可以敵大，寡固不可以敵衆，弱固不可以敵彊。海內之地方千里者九，齊集有其一。以一服八，何以異於鄒敵楚哉？」\jiazhu[1]{固，辭也。言小弱固不如彊大。集會齊地可方千里，譬一州耳。今欲以一州服八州，猶鄒欲敵楚。}

「蓋亦反其本矣！」\jiazhu[1]{王欲服之之道，蓋當反王道之本。}

「今王發政施仁，使天下仕者皆欲立於王之朝，耕者皆欲耕於王之野，商賈皆欲藏於王之市，行旅皆欲出於王之塗，天下之欲疾其君者皆欲赴愬於王。其若是，孰能禦之？」\jiazhu[1]{反本道，行仁政。若此則天下歸之，誰能止之者。}

王曰：「吾惛，不能進於是矣。願夫子輔吾志，明以敎我。我雖不敏，請甞試之。」\jiazhu[1]{王言我情思惛亂，不能進行此仁政，不知所當施行也。欲使孟子明言其道以敎訓之，我雖不敏，願甞使少行之也。}

曰：「無恒産而有恒心者，惟士爲能。若民，則無恒産，因無恒心。」\jiazhu[1]{孟子爲王陳其法也。恒，常也。産，生也。恒産則民常可以生之業也。恒心，人常有所善心也。惟有學士之心者，雖窮不失道，不求苟得耳。凡民迫於飢寒，則不能守其常善之心。}

「苟無恒心，放辟邪侈，無不爲已。及陷於罪，然後從而刑之，是罔民也。」\jiazhu[1]{民誠無恒心，放溢辟邪，侈於姧利，犯罪觸刑，無所不爲。乃就刑之，是由張羅罔以罔民者也。}

「焉有仁人在位罔民而可爲也？」\jiazhu[1]{安有仁人爲君罔陷其民，是政何可爲也。}

「是故明君制民之産，必使仰足以事父母，俯足以畜妻子，樂歲終身飽，凶年免於死亡。然後驅而之善，故民之從之也輕。」\jiazhu[1]{言衣食足知榮辱，故民從之敎化輕易也。}

「今也制民之産，仰不足以事父母，俯不足以畜妻子，樂歲終身苦，凶年不免於死亡。此惟救死而恐不贍，奚暇治禮義哉？」\jiazhu[1]{言今民困窮，救死恐凍餓而不給，何暇脩禮行義也。}

「王欲行之，則盍反其本矣！」\jiazhu[1]{王欲行之，則盍反其本矣！}

「五畒之宅，樹之以桑，五十者可以衣帛矣。雞豚狗彘之畜，無失其時，七十者可以食肉矣。百畒之田，勿奪其時，八口之家可以無飢矣。謹庠序之敎，申之以孝悌之義，頒白者不負戴於道路矣。老者衣帛食肉，黎民不飢不寒，然而不王者，未之有也。」\jiazhu[1]{其說與上同。八口之家，次上農夫也。孟子所以重言此者，此乃王政之本，常生之道。故爲齊梁之君各具陳之，當章究義，不嫌其重也。}\par \jiazhu[1]{章指言：典籍攸載，帝王道純，桓文之事譎正相紛，撥亂反正，聖意弗珍，故曰後世無傳未聞。仁不施人，猶不成德。釁鍾易牲，民不被澤，王請甞試，欲踐其路，荅以反本，惟是爲要。此蓋孟子不屈道之言也。}
